% 考古天文学（综述）
% license CCBYNCSA3
% type Wiki

本文根据 CC-BY-SA 协议转载翻译自维基百科\href{https://en.wikipedia.org/wiki/Archaeoastronomy}{相关文章}。

\begin{figure}[ht]
\centering
\includegraphics[width=6cm]{./figures/b9737bc319a7b6fe.png}
\caption{} \label{fig_ArmY_1}
\end{figure}
考古天文学（Archaeoastronomy，又作 Archeoastronomy）是一门跨学科（interdisciplinary）[1]或多学科（multidisciplinary）[2]的研究领域，主要探讨古代人类如何“理解天空中的现象、如何利用这些现象，以及天空在其文化中所扮演的角色”[3]。Clive Ruggles 指出，将考古天文学简单地视为“古代天文学的研究”是误导性的，因为现代天文学是一门科学学科，而考古天文学则关注不同文化对天象的象征性与文化性解释[4][5]。考古天文学常与民族天文学（ethnoastronomy）并行研究，后者是一门人类学分支，专门研究当代社会中的观星活动与文化观念。此外，考古天文学还与历史天文学（historical astronomy）密切相关—— 后者利用古代的天象记录来解决天文学问题，以及天文学史（history of astronomy），即通过书面文献来评估古代天文学实践的发展历程[6]。


考古天文学（Archaeoastronomy）使用多种方法来揭示古代天文实践的证据，这些方法包括考古学、 人类学、 天文学、 统计学与概率论、以及历史学[7]。由于这些方法来源多样、数据类型差异巨大，如何将它们整合为一致而连贯的论证，一直是考古天文学家面临的长期挑战[8]。考古天文学在景观考古学（landscape archaeology）和认知考古学（cognitive archaeology）中占据互补地位。物质证据及其与天空的关联可以揭示更广阔的自然景观如何融入信仰体系，例如玛雅天文学（Mayan astronomy）与农业周期之间的关系[9]。其他结合认知与景观概念的研究实例包括：探讨定居点道路布局中所体现的宇宙秩序（cosmic order）[10][11]。  


考古天文学可适用于所有文化与各个历史时期。  
尽管不同文化对天空的意义理解差异显著，  
但在研究古代信仰时，仍然可以跨文化地应用一套科学方法[12]。  
或许正是因为必须在\textbf{社会学维度}与\textbf{科学方法}之间取得平衡，  
Clive Ruggles 才将考古天文学形容为——  
“一个在一端产出高质量学术研究，  
而在另一端则存在近乎疯狂的无拘臆测的领域”[13]。
```
